% Gerado por ferramentas/gerar.py; não edite à mão.
\section*{Exercícios}\markright{Exercícios}
\addcontentsline{toc}{section}{Exercícios}
Os pontos antes de cada número indicam a dificuldade: \nivel{1} fácil, \nivel{2} médio, \nivel{3} difícil e \nivel{4} desafio. Os exercícios marcados com \estrela{} têm solução completa, com animação do gráfico, no repositório do livro: \url{https://github.com/AleffAzevedo/calculo-1-pt/blob/main/exercicios/cap5/README.md}. As respostas estão no fim do capítulo.

\begin{itemize}[leftmargin=5.2em,labelsep=0.6em,itemsep=4pt]
\item[\exnum{1}{5.1}] Calcule $\displaystyle\sum_{i=1}^{20}(3i-2)$ e $\displaystyle\sum_{i=1}^{10}i^2$ usando as fórmulas de somatório.
\item[\exnum{1}{5.2}] Calcule $L_4$ e $R_4$ para $f(x)=x^2+1$ em $[0,2]$. Qual subestima a área? Por quê?
\item[\exnum{1}{5.3}] Calcule a soma do ponto médio $M_4$ para $f(x)=\dfrac1x$ em $[1,3]$.
\item[\exnum{1}{5.4}] Interprete $\int_0^3\sqrt{9-x^2}\,dx$ como área e calcule-a sem antiderivadas.
\item[\exnum{1}{5.5}] Use áreas de figuras geométricas para calcular $\int_{-2}^{4}|x-1|\,dx$.
\item[\exnum{1}{5.6}] Sabendo que $\int_0^5f=7$, $\int_0^3f=4$ e $\int_0^5g=-2$, calcule $\int_3^5f$, $\int_5^0(2f-3g)$ e $\int_3^3g$.
\item[\exnum{1}{5.7}] Calcule o valor médio de $f(x)=x^2$ em $[0,3]$ e encontre $c$ com $f(c)=f_{\text{méd}}$.
\item[\exnum{1}{5.8}] Calcule $\dfrac{d}{dx}\displaystyle\int_2^x\frac{\sin t}{t}\,dt$ e $\dfrac{d}{dx}\displaystyle\int_0^{x^3}\sqrt{1+t^2}\,dt$.
\item[\exnum{1}{5.9}] Calcule $\displaystyle\int_1^2\Big(3x^2-\frac2{x^3}\Big)dx$.
\item[\exnum{1}{5.10}] Calcule $\displaystyle\int_0^{\pi/3}\sec^2x\,dx$ e $\displaystyle\int_0^{\pi}(\cos x+2)\,dx$.
\item[\exnum{1}{5.11}] Calcule $\displaystyle\int_1^4\frac{x-1}{\sqrt x}\,dx$.
\item[\exnum{1}{5.12}] Calcule $\displaystyle\int_0^1(e^x-2^x)\,dx$.
\item[\exnum{1}{5.13}] Calcule $\displaystyle\int(2x+1)^7\,dx$ e $\displaystyle\int\cos(3x-1)\,dx$.
\item[\exnum{1}{5.14}] Calcule $\displaystyle\int x\sqrt{x^2+1}\,dx$.
\item[\exnum{1}{5.15}] Calcule $\displaystyle\int\frac{x}{x^2+4}\,dx$ e $\displaystyle\int\frac{dx}{2x+5}$.
\item[\exnum{1}{5.16}] Calcule $\displaystyle\int e^{5x}\,dx$ e $\displaystyle\int x^2e^{x^3}\,dx$.
\item[\exnum{1}{5.17}] Calcule $\displaystyle\int\frac{dx}{\sqrt{9-x^2}}$ e $\displaystyle\int\frac{dx}{16+x^2}$.
\item[\exnum{1}{5.18}] Calcule $\displaystyle\int_0^{\pi/2}\sin x\cos^3x\,dx$.
\item[\exnum{1}{5.19}] Calcule $\displaystyle\int_{-1}^{1}(x^5-3x^3+x^2)\,dx$ usando simetria.
\item[\exnum{1}{5.20}] Uma partícula tem $v(t)=2t-4$ m/s em $[0,5]$. Calcule o deslocamento e a distância percorrida.
\item[\exnum{2}{5.21}] \estrela\ \textbf{Somas de Riemann convergindo.} Para $f(x)=4-x^2$ em $[0,2]$, calcule $L_4$ e $R_4$, escreva $R_n$ em função de $n$ e
encontre a área exata como $\lim_{n\to\infty}R_n$.
\item[\exnum{2}{5.22}] \estrela\ \textbf{Integral pela definição.} Use a definição (limite de somas de Riemann com extremidades direitas) para calcular
$\int_0^1x^3\,dx$.
\item[\exnum{2}{5.23}] Escreva $\displaystyle\lim_{n\to\infty}\sum_{i=1}^n\Big(1+\frac{2i}n\Big)^2\frac2n$ como integral definida e calcule.
\item[\exnum{2}{5.24}] Mostre, sem calcular, que $2\le\displaystyle\int_0^2\sqrt{1+x^3}\,dx\le6$.
\item[\exnum{2}{5.25}] \estrela\ \textbf{A função área.} Seja $g(x)=\displaystyle\int_0^x(t-1)(t-3)\,dt$ para $0\le x\le4$. Use o TFC para encontrar
$g'(x)$, os intervalos em que $g$ cresce e decresce, e os extremos locais de $g$. Calcule
$g(1)$, $g(3)$ e $g(4)$.
\item[\exnum{2}{5.26}] \estrela\ \textbf{TFC com limites variáveis.} Calcule $\dfrac{d}{dx}\displaystyle\int_x^{x^2}e^{t^2}\,dt$.
\item[\exnum{2}{5.27}] Calcule $\dfrac{d}{dx}\displaystyle\int_{\cos x}^{1}\frac{dt}{1+t^2}$.
\item[\exnum{2}{5.28}] Calcule $\displaystyle\int_0^3|x^2-4|\,dx$.
\item[\exnum{2}{5.29}] \estrela\ \textbf{Deslocamento e distância.} Uma partícula tem velocidade $v(t)=3t^2-12t+9$ m/s para $0\le t\le4$. Calcule o
deslocamento e a distância percorrida. Se $s(0)=2$, onde ela está em $t=4$?
\item[\exnum{2}{5.30}] \estrela\ \textbf{Valor médio de uma função.} Calcule o valor médio de $f(x)=\sin x$ em $[0,\pi]$ e encontre todos os $c$ garantidos pelo
TVM para integrais.
\item[\exnum{2}{5.31}] Calcule $\displaystyle\int x\sqrt{x+2}\,dx$.
\item[\exnum{2}{5.32}] Calcule $\displaystyle\int\frac{\cos x}{\sin^2x}\,dx$ e $\displaystyle\int\sec^2x\,\tan^3x\,dx$.
\item[\exnum{2}{5.33}] \estrela\ \textbf{Substituição em integral definida.} Calcule $\displaystyle\int_0^2x\sqrt{4-x^2}\,dx$ trocando os limites de integração, e interprete
a substituição geometricamente.
\item[\exnum{2}{5.34}] Calcule $\displaystyle\int_0^{1}\frac{e^x}{1+e^x}\,dx$.
\item[\exnum{2}{5.35}] Calcule $\displaystyle\int\frac{(\ln x)^2}{x}\,dx$ e $\displaystyle\int\frac{dx}{x\ln x}$.
\item[\exnum{2}{5.36}] Calcule $\displaystyle\int\frac{e^{x}}{1+e^{2x}}\,dx$.
\item[\exnum{2}{5.37}] Calcule $\displaystyle\int_0^{1/2}\frac{dx}{\sqrt{1-x^2}}$ e $\displaystyle\int_{0}^{\sqrt3}\frac{dx}{1+x^2}$.
\item[\exnum{2}{5.38}] Calcule $\displaystyle\int\frac{dx}{x^2-6x+13}$ completando o quadrado.
\item[\exnum{2}{5.39}] \estrela\ \textbf{Simetria salva contas.} Calcule $\displaystyle\int_{-\pi}^{\pi}\big(x^3\cos x+2\cos x+x\big)\,dx$ usando simetria.
\item[\exnum{2}{5.40}] Calcule $\displaystyle\int\frac{x+3}{x^2+1}\,dx$ (separe em duas integrais).
\item[\exnum{3}{5.41}] \estrela\ \textbf{Calculando $\pi$ com uma integral.} Mostre que $\displaystyle\int_0^1\frac{4}{1+x^2}\,dx=\pi$ e use a regra do ponto médio com $n=4$ para
estimar $\pi$.
\item[\exnum{3}{5.42}] \estrela\ \textbf{Vazão e volume.} Água entra num tanque à taxa $r(t)=200-4t$ L/min, $0\le t\le50$, que inicialmente tem
1000 L. Quanto há no tanque em $t=50$? Em que instante o volume cresce mais rapidamente?
Esboce $V(t)$.
\item[\exnum{3}{5.43}] \estrela\ \textbf{Integral da tangente e da secante.} Mostre que $\int\tan x\,dx=\ln|\sec x|+C$ e que $\int\sec x\,dx=\ln|\sec x+\tan x|+C$.
Calcule $\int_0^{\pi/4}\sec x\,dx$.
\item[\exnum{3}{5.44}] Use somas de Riemann (extremidades direitas) para calcular $\int_1^3(2x+1)\,dx$ pela definição.
\item[\exnum{3}{5.45}] Encontre $f$ contínua e a constante $a$ tais que $\displaystyle6+\int_a^x\frac{f(t)}{t^2}\,dt=2\sqrt x$ para $x>0$.
\item[\exnum{3}{5.46}] Encontre os extremos locais e os pontos de inflexão de $F(x)=\displaystyle\int_0^x\frac{t-2}{1+t^2}\,dt$, $x\ge0$.
\item[\exnum{3}{5.47}] Calcule $\displaystyle\lim_{x\to0}\frac1{x^3}\int_0^x\sin(t^2)\,dt$.
\item[\exnum{3}{5.48}] Calcule $\displaystyle\int_0^{\pi}|\cos x|\,dx$ e $\displaystyle\int_0^{2\pi}|\sin x-\cos x|\,dx$.
\item[\exnum{3}{5.49}] A taxa de vendas de um produto é $r(t)=100\,t\,e^{-t^2/10}$ unidades/semana. Quantas unidades são vendidas nas 4 primeiras semanas? E no longo prazo?
\item[\exnum{3}{5.50}] Calcule $\displaystyle\int\frac{x^3}{\sqrt{1-x^2}}\,dx$ (dica: $u=1-x^2$).
\item[\exnum{3}{5.51}] Calcule $\displaystyle\int\frac{dx}{\sqrt{e^{2x}-1}}$ (dica: $u=e^x$).
\item[\exnum{3}{5.52}] Calcule $\displaystyle\int\frac{\sin x}{1+\cos^2x}\,dx$ e $\displaystyle\int\frac{dx}{x\sqrt{x^2-4}}$.
\item[\exnum{3}{5.53}] Calcule $\displaystyle\int\frac{dx}{\sqrt{-x^2+4x}}$ completando o quadrado.
\item[\exnum{3}{5.54}] Mostre que, se $f$ é contínua e periódica de período $T$, então $\int_a^{a+T}f$ não depende de $a$.
\item[\exnum{3}{5.55}] Um carro freia com $a(t)=-8$ m/s² a partir de $30$ m/s. Use integrais para achar a distância até parar.
\item[\exnum{3}{5.56}] O valor eficaz (RMS) de $f$ em $[a,b]$ é $\sqrt{\tfrac1{b-a}\int_a^bf^2}$. Calcule o RMS de $f(t)=\sin t$ em $[0,2\pi]$. (Dica: $\sin^2t=\tfrac{1-\cos2t}2$.)
\item[\exnum{3}{5.57}] Seja $g(x)=\int_0^x f(t)\,dt$, com $f$ crescente e $f(0)=0$. Mostre que $g$ é côncava para cima e que $g(x)\ge0$ para $x\ge0$.
\item[\exnum{3}{5.58}] Calcule $\displaystyle\lim_{n\to\infty}\frac1n\sum_{i=1}^n\sqrt{\frac in}$ interpretando como integral.
\item[\exnum{3}{5.59}] Calcule $\displaystyle\lim_{n\to\infty}\sum_{i=1}^n\frac{n}{n^2+i^2}$.
\item[\exnum{3}{5.60}] Calcule $\displaystyle\int_0^{\ln2}e^x\sqrt{e^x-1}\,dx$ e $\displaystyle\int_1^{e^2}\frac{dx}{x\sqrt{\ln x}}$.
\item[\exnum{4}{5.61}] Prove que $\displaystyle\int_0^{\pi/2}\frac{\sin^nx}{\sin^nx+\cos^nx}\,dx=\frac\pi4$ para todo $n>0$. (Dica: $u=\tfrac\pi2-x$.)
\item[\exnum{4}{5.62}] Mostre que $\ln x=\int_1^x\frac{dt}{t}$ satisfaz $\ln(ab)=\ln a+\ln b$ usando apenas propriedades da integral e substituição.
\item[\exnum{4}{5.63}] Calcule $\displaystyle\int_0^1\frac{\ln(1+x)}{1+x^2}\,dx$. (Dica: $x=\tan\theta$ e depois $\theta\mapsto\tfrac\pi4-\theta$.)
\item[\exnum{4}{5.64}] Encontre todas as funções contínuas $f$ com $\displaystyle\int_0^xf(t)\,dt=x^2+\int_x^1t^2f(t)\,dt$ para todo $x$.
\item[\exnum{4}{5.65}] Prove a desigualdade $\ln(1+n)<1+\tfrac12+\dots+\tfrac1n<1+\ln n$ comparando somas com $\int\frac{dx}x$.
\end{itemize}
\subsection*{Respostas}\markright{Respostas}
{\small\begin{multicols}{2}\raggedright
\begin{itemize}[leftmargin=2.7em,labelsep=0.5em,itemsep=1pt]
\item[\textbf{5.1}] $590$ e $385$.
\item[\textbf{5.2}] $L_4=3{,}75$; $R_4=5{,}75$; $L_4$ subestima, pois $f$ é crescente.
\item[\textbf{5.3}] $M_4=\tfrac{2}{5}+\tfrac27+\tfrac29+\tfrac2{11}\approx1{,}0898$ (exato: $\ln3\approx1{,}0986$).
\item[\textbf{5.4}] Um quarto do círculo de raio 3: $\tfrac{9\pi}4$.
\item[\textbf{5.5}] $\tfrac92+\tfrac92=9$.
\item[\textbf{5.6}] $3$; $-20$; $0$.
\item[\textbf{5.7}] $3$; $c=\sqrt3$.
\item[\textbf{5.8}] $\tfrac{\sin x}x$; $3x^2\sqrt{1+x^6}$.
\item[\textbf{5.9}] $\tfrac{25}4$.
\item[\textbf{5.10}] $\sqrt3$; $2\pi$.
\item[\textbf{5.11}] $\tfrac83$.
\item[\textbf{5.12}] $e-1-\tfrac1{\ln2}\approx0{,}276$.
\item[\textbf{5.13}] $\tfrac{(2x+1)^8}{16}+C$; $\tfrac13\sin(3x-1)+C$.
\item[\textbf{5.14}] $\tfrac13(x^2+1)^{3/2}+C$.
\item[\textbf{5.15}] $\tfrac12\ln(x^2+4)+C$; $\tfrac12\ln|2x+5|+C$.
\item[\textbf{5.16}] $\tfrac15e^{5x}+C$; $\tfrac13e^{x^3}+C$.
\item[\textbf{5.17}] $\sin^{-1}\tfrac x3+C$; $\tfrac14\tan^{-1}\tfrac x4+C$.
\item[\textbf{5.18}] $\tfrac14$.
\item[\textbf{5.19}] $\tfrac23$.
\item[\textbf{5.20}] $5$ m e $13$ m.
\item[\textbf{5.21}] $L_4=6{,}25$; $R_4=4{,}25$; $R_n=8-\tfrac{4(n+1)(2n+1)}{3n^2}$; área $\tfrac{16}3$.
\item[\textbf{5.22}] $\tfrac14$.
\item[\textbf{5.23}] $\int_1^3x^2\,dx=\tfrac{26}3$.
\item[\textbf{5.24}] $1\le\sqrt{1+x^3}\le3$ em $[0,2]$.
\item[\textbf{5.25}] $g'(x)=(x-1)(x-3)$; máx. local $g(1)=\tfrac43$; mín. local $g(3)=0$; $g(4)=\tfrac43$.
\item[\textbf{5.26}] $2x\,e^{x^4}-e^{x^2}$.
\item[\textbf{5.27}] $\dfrac{\sin x}{1+\cos^2x}$.
\item[\textbf{5.28}] $\tfrac{23}3$.
\item[\textbf{5.29}] Deslocamento 4 m; distância 12 m; $s(4)=6$.
\item[\textbf{5.30}] $f_{\text{méd}}=\tfrac2\pi\approx0{,}637$; $c=\sin^{-1}\tfrac2\pi\approx0{,}690$ ou $c=\pi-0{,}690\approx2{,}451$.
\item[\textbf{5.31}] $\tfrac25(x+2)^{5/2}-\tfrac43(x+2)^{3/2}+C$.
\item[\textbf{5.32}] $-\csc x+C$; $\tfrac14\tan^4x+C$.
\item[\textbf{5.33}] $\tfrac83$.
\item[\textbf{5.34}] $\ln\tfrac{1+e}2\approx0{,}620$.
\item[\textbf{5.35}] $\tfrac13(\ln x)^3+C$; $\ln|\ln x|+C$.
\item[\textbf{5.36}] $\tan^{-1}(e^x)+C$.
\item[\textbf{5.37}] $\tfrac\pi6$; $\tfrac\pi3$.
\item[\textbf{5.38}] $\tfrac12\tan^{-1}\tfrac{x-3}2+C$.
\item[\textbf{5.39}] $0$.
\item[\textbf{5.40}] $\tfrac12\ln(x^2+1)+3\tan^{-1}x+C$.
\item[\textbf{5.41}] $4\tan^{-1}1=\pi$; $M_4\approx3{,}14680$.
\item[\textbf{5.42}] $6000$ L; em $t=0$; $V(t)=1000+200t-2t^2$, parábola crescente até o vértice em $t=50$.
\item[\textbf{5.43}] $\ln(\sqrt2+1)\approx0{,}8814$.
\item[\textbf{5.44}] $10$.
\item[\textbf{5.45}] $f(x)=x^{3/2}$, $a=9$.
\item[\textbf{5.46}] Mín. local em $x=2$; inflexão em $x=2+\sqrt5$.
\item[\textbf{5.47}] $\tfrac13$ (L'Hôpital + TFC).
\item[\textbf{5.48}] $2$; $4\sqrt2$.
\item[\textbf{5.49}] $500(1-e^{-1{,}6})\approx399$; tende a $500$.
\item[\textbf{5.50}] $-\sqrt{1-x^2}+\tfrac13(1-x^2)^{3/2}+C$.
\item[\textbf{5.51}] $\sec^{-1}(e^x)+C$.
\item[\textbf{5.52}] $-\tan^{-1}(\cos x)+C$; $\tfrac12\sec^{-1}\tfrac{|x|}2+C$.
\item[\textbf{5.53}] $\sin^{-1}\tfrac{x-2}2+C$.
\item[\textbf{5.54}] Derive $G(a)=\int_a^{a+T}f$: $G'(a)=f(a+T)-f(a)=0$.
\item[\textbf{5.55}] $56{,}25$ m.
\item[\textbf{5.56}] $\tfrac1{\sqrt2}\approx0{,}707$.
\item[\textbf{5.57}] $g''=f'\ge0$; $g$ acumula valores $f(t)\ge0$.
\item[\textbf{5.58}] $\int_0^1\sqrt x\,dx=\tfrac23$.
\item[\textbf{5.59}] $\int_0^1\tfrac{dx}{1+x^2}=\tfrac\pi4$.
\item[\textbf{5.60}] $\tfrac23$; $2\sqrt2$.
\item[\textbf{5.61}] A integral $I$ satisfaz $I=\int_0^{\pi/2}\tfrac{\cos^n}{\sin^n+\cos^n}$; some as duas: $2I=\tfrac\pi2$.
\item[\textbf{5.62}] $\int_a^{ab}\tfrac{dt}t=\int_1^b\tfrac{du}u$ com $t=au$.
\item[\textbf{5.63}] $\tfrac\pi8\ln2$.
\item[\textbf{5.64}] Nenhuma: derivando, $f=\tfrac{2x}{1+x^2}$, mas em $x=0$ seria preciso $\int_0^1t^2f=0$, e essa integral vale $1-\ln2$.
\item[\textbf{5.65}] Retângulos acima e abaixo de $y=1/x$ em $[1,n+1]$ e $[1,n]$.
\end{itemize}\end{multicols}}
